\chapter{Fitting models to data}
\label{ch:fitting}

The module \code{roadrunner\_nitro.fit} estimates model parameters from measured data. It has two levels:

\begin{itemize}
  \item \code{fit.fit\_timecourse(model, data, params)}: one time course, in one line, with confidence intervals
        and a plot (section~\ref{sec:timecourse}).
  \item \code{fit.Problem}: several experiments, observables that are formulas of the model, measurement error
        models, pre-equilibration, PEtab files and every option of the optimizers (section~\ref{sec:problem}).
\end{itemize}

The short form builds a \code{Problem} internally, so you can move from one to the other without starting over.

\section{The short form: one time course}
\label{sec:timecourse}

We measured \code{S1} at ten times and want the two rate constants:

\begin{lstlisting}[language=Python]
from roadrunner_nitro import fit

model = """
    J1: $X0 -> S1; k1*X0
    J2: S1 -> ;    k2*S1
    X0 = 10; S1 = 0; k1 = 1; k2 = 1
"""
data = {'time': [0, 2, 4, 6, 8, 10, 15, 20, 25, 30],
        'S1':   [0.0, 7.0, 14.0, 17.2, 19.7, 21.5, 22.7, 24.4, 24.4, 26.6]}

result = fit.fit_timecourse(model, data, ['k1', 'k2'])
print(result)
result.plot()                  # the data and the fitted curve (needs matplotlib)
\end{lstlisting}
\begin{lstlisting}[style=out]
Fit of 2 parameters to 10 data points (9 of 10 starts reached the best fit)
  k1 = 0.47004    (95% interval 0.4275 to 0.5168)
  k2 = 0.18585    (95% interval 0.1652 to 0.2091)
  measurement error estimated from the scatter of the data: 0.7813
\end{lstlisting}

\paragraph{The model} can be Antimony text, SBML text, the name of a model file, or a \code{RoadRunner} object.

\paragraph{The data} is a dictionary, a pandas table or a \code{simulate()} result, with a \code{'time'} column
and one column per measured quantity, named by its model symbol: a species, \code{'[S1]'} for its
concentration, a parameter, or an assignment-rule variable. Missing values (NaN) are skipped. The simulation
starts at time 0 from the model's initial state.

\paragraph{The parameters} are a list of ids, each searched from 1/1000 to 1000 times its value in the model on
a logarithmic scale, or a dictionary giving ranges, \code{\{'k1': (0.01, 10)\}}, optionally with a starting
value, \code{\{'k1': (0.01, 10, 0.5)\}}. The model's values are the first starting point.

\paragraph{The measurement error} is estimated from the scatter of the data around the fit, as
\code{scipy.optimize.curve\_fit} and lmfit do, unless you give it: \code{sigma=0.5} for every column, or a
dictionary per column. The confidence intervals use it.

\paragraph{Options.} \code{n\_starts} (default 10) is the number of starting points of the multistart fit, and
\code{seed} the random seed that places them. \code{global\_method} searches the whole range first with
\code{'pso'}, \code{'de'} or \code{'shade'} (section~\ref{sec:optimizers}); \code{n\_starts} is then the number of
the search's best points refined. \code{level} (default 0.95) is the confidence level of the intervals.

\paragraph{The result} has \code{values} and \code{intervals} (dictionaries by id), \code{identifiable},
\code{sigma}, \code{simulate()} and \code{plot()}, and the full \code{problem}, \code{result} and \code{fisher}
objects behind them.

\section{Fitting problems}
\label{sec:problem}

A fitting problem has five parts, the same as a PEtab problem:

\begin{center}
\begin{tabularx}{\textwidth}{lY}
\toprule
Part & What it says \\
\midrule
Model & the SBML model \\
Parameters & which parameters are estimated, their bounds and scales; fixed values for the others \\
Conditions & the experiments: which parameters or initial values each one changes \\
Observables & what was measured, as a formula of the model, and its measurement error \\
Measurements & the data: condition, observable, time, value \\
\bottomrule
\end{tabularx}
\end{center}

\subsection{Building a problem}

\begin{lstlisting}[language=Python]
p = fit.Problem(sbml)            # or fit.Problem.from_antimony(text), .from_roadrunner(r)

p.add_parameter('k1', lower=1e-3, upper=1e3)   # estimated, within the bounds
p.add_parameter('k2', lower=1e-3, upper=1e3)

p.add_condition('normal')                      # the model as it is
p.add_condition('half', {'X0': 5})             # an experiment with X0 = 5

p.add_observable('S1_obs', 'S1', noise_formula='0.8')
p.add_measurements('normal', 'S1_obs', times, values)
p.add_measurements('half', 'S1_obs', times, values_half)
print(p)
\end{lstlisting}
\begin{lstlisting}[style=out]
<fit.Problem: 2 free parameters, 20 measurements, 2 conditions, 1 observables>
\end{lstlisting}

\paragraph{Parameters.} \code{add\_parameter(id, scale, lower, upper, nominal, estimate)}. The
\code{scale} is the one the optimizer works on: \code{'lin'}, \code{'log'} or \code{'log10'}. Left out, a
positive range spanning a factor of 100 or more is searched on a log10 scale, otherwise linearly; rate constants
are almost always best searched on a log scale. A parameter with \code{estimate=False} is fixed at its
\code{nominal} value. An id that is not in the model is an extra parameter that observables can use, such as a
scale factor or an offset.

\paragraph{Conditions.} \code{add\_condition(id, overrides)} gives the values an experiment changes: parameters,
initial values of species (concentrations), or compartment sizes. A value can be a number or the id of a
parameter in the parameter table, so an initial value can itself be estimated.

\paragraph{Observables.} \code{add\_observable(id, formula, noise\_formula, transformation='lin',
noise='normal')}. The formula uses model symbols, parameter ids and \code{time}, so observables such as
\code{'scale * (S1 + S2)'} are written directly. The noise formula is the standard deviation of the measurement
error; it can depend on parameters (an estimated error) or on the model state (a relative error such as
\code{'0.1 * S1'}). \code{transformation} is \code{'lin'}, \code{'log'} or \code{'log10'} (the data are
compared on that scale, so log-normal errors are \code{'log'} with normal noise), and \code{noise} is
\code{'normal'} or \code{'laplace'}.

\paragraph{Measurements.} \code{add\_measurements(condition, observable, times, values)}. A time of
\code{numpy.inf} is a measurement at steady state.

\subsection{Pre-equilibration and steady states}

Many experiments start from a system at rest: cells are kept unstimulated until they settle, then treated. A
\emph{pre-equilibration condition} says so:

\begin{lstlisting}[language=Python]
p.add_condition('unstimulated', {'ligand': 0})
p.add_condition('stimulated', {'ligand': 10})
p.add_measurements('stimulated', 'pERK', times, values,
                   preequilibration='unstimulated')
\end{lstlisting}

The model is brought to its steady state with the values of \code{'unstimulated'}, and the simulation of
\code{'stimulated'} starts from that state. States that \code{'stimulated'} itself sets keep the values it sets.
The steady state is searched from the model's initial state at every evaluation, so the result never depends on
earlier evaluations.

A measurement at time \code{inf} is taken at the steady state reached from the last finite measurement time. Both
work with conservation laws: a conserved total keeps the value it has where the steady-state search starts, so
the steady state depends on the initial state as it should, and the gradient accounts for it.

\subsection{PEtab problems}

PEtab is a standard file format for fitting problems (\url{https://petab.readthedocs.io}). A PEtab version 1
problem is read from its YAML file:

\begin{lstlisting}[language=Python]
p = fit.Problem.from_petab('problem.yaml')
\end{lstlisting}

Supported: the parameter, condition, observable and measurement tables; parameter scales and observable
transformations; normal and Laplace noise; observable and noise parameter placeholders
(\code{observableParameters}, \code{noiseParameters}); pre-equilibration; steady-state measurements; condition
values for parameters, initial values and compartment sizes. Features that are not supported yet (priors, models
with events) are refused with an error naming the table, row and column, never silently ignored.

\section{The objective and its gradient}

The objective is the negative log-likelihood of the data, as in PEtab, including its constant terms. For a
measurement $m$ with simulated observable $y$, transformation $h$ and standard deviation $\sigma$, normal noise
contributes
\[
  \tfrac12 \log\left(2\pi\sigma^2\right) + \tfrac12\left(\frac{h(y) - h(m)}{\sigma}\right)^2
\]
and Laplace noise $\log(2\sigma) + |h(y) - h(m)|/\sigma$. When the noise is normal and depends on no estimated
parameter and no state, the problem is a least-squares problem (\code{p.least\_squares}), and the residuals
$(h(y)-h(m))/\sigma$ and their Jacobian are available too.

Gradients are exact: the sensitivities of the model to every estimated parameter are integrated with the model,
sharing its steps, rather than approximated by finite differences. They include the effect of parameters on
initial values (through initial assignments), on switch times in the model, and on steady states.

\begin{lstlisting}[language=Python]
theta = p.nominal                  # parameters on their scales (here log10)
p.objective(theta)
f, g = p.gradient(theta)           # objective and its gradient
r, J = p.residual_jacobian(theta)  # least-squares problems only
p.simulate(theta)                  # the simulated observable of every measurement
p.to_linear(theta)                 # {'k1': ..., 'k2': ...} on the linear scale
\end{lstlisting}

The integrator tolerances of every simulation are \code{p.rtol} and \code{p.atol} (default $10^{-8}$ and
$10^{-12}$).

\section{Optimizers}
\label{sec:optimizers}

\code{p.fit()} minimizes the objective within the bounds. It has two ways to search.

\paragraph{Multistart (the default).} \code{n\_starts} starting points are spread over the box of the bounds by
a Latin hypercube sample, and a local optimizer runs from each, in parallel. The local optimizer is chosen
automatically: trust-region reflective least squares for least-squares problems, using the residuals and their
exact Jacobian, and L-BFGS-B for other problems, using the objective and its exact gradient.

\begin{lstlisting}[language=Python]
r = p.fit(n_starts=20)
print(r)
\end{lstlisting}
\begin{lstlisting}[style=out]
<fit.FitResult: best objective 18.68537727, 18 of 20 starts within 1e-3 of it, 20 converged>
  k1 = 0.47004482
  k2 = 0.18585473
\end{lstlisting}

That most starts reach the same best value is the usual sign that it is the global optimum.

\paragraph{A global search, then local refinement.} With \code{global\_method}, a population-based method first
searches the whole box using the objective only, and its \code{n\_starts} best distinct points are then refined
by the local optimizer:

\begin{center}
\begin{tabularx}{\textwidth}{lY}
\toprule
\code{global\_method} & Method \\
\midrule
\code{'pso'} & particle swarm optimization, SPSO-2011 (Clerc 2011; Zambrano-Bigiarini, Clerc and Rojas 2013):
  40 particles, adaptive random topology \\
\code{'de'} & differential evolution, DE/rand/1/bin (Storn and Price 1997): $F = 0.5$, $CR = 0.9$, population
  10 per parameter between 20 and 200 \\
\code{'shade'} & SHADE (Tanabe and Fukunaga 2013), differential evolution that adapts its own $F$ and $CR$:
  population 100 \\
\bottomrule
\end{tabularx}
\end{center}

\begin{lstlisting}[language=Python]
r = p.fit(global_method='de', n_starts=3)          # refine the 3 best points found
r = p.fit(global_method='pso', max_evals=20000)    # a larger search budget
r = p.fit(global_method='shade', polish=False)     # the search's points, unrefined
r.global_stage    # method, best, nfev, generations, message of the search
\end{lstlisting}

The search stops after \code{max\_evals} objective evaluations (default 1000 per parameter, at least 10\,000),
when the values of its members agree closely, or when it stops improving. \code{population} changes the number
of members. No method is best on every problem; model fitting rewards trying more than one.

\paragraph{Common options.} \code{seed} chooses every random number: the same seed gives the same result, to the
last bit, whatever the number of threads. \code{threads} limits the cores used (default: all). \code{x0} adds a
starting point of your own. \code{callback(finished, total, best)} is called after each start (during a global
search, after each generation with the number of evaluations); returning \code{True} cancels the fit, and starts
already running keep their best point so far. \code{ftol}, \code{xtol}, \code{gtol} and \code{max\_nfev} set the
local optimizer's stopping rules.

\subsection{Results}

\code{fit} returns a \code{FitResult} with every start, best first: \code{best} (the parameters on their scales),
\code{best\_value}, \code{best\_linear} (a dictionary of linear values), \code{x}, \code{x0} and \code{value} for
every start, \code{converged}, \code{message}, evaluation counts, and \code{n\_at\_best()}, the number of starts
that reached the best value. \code{r.to\_dataframe()} gives a pandas table with one row per start, and
\code{r.waterfall()} the sorted values for a waterfall plot.

\section{Confidence intervals}

\code{p.fisher(theta, level=0.95)} gives Wald intervals from the Fisher information, the Gauss-Newton
approximation $J^\mathsf{T} J$ of the Hessian, at \code{theta}:

\begin{lstlisting}[language=Python]
print(p.fisher(r.best))
\end{lstlisting}
\begin{lstlisting}[style=out]
<fit.FisherResult: 95% intervals>
  k1: [0.430942, 0.512696]
  k2: [0.166829, 0.207051]
\end{lstlisting}

The intervals are symmetric on the parameter's scale, so on a log scale they are asymmetric on the linear scale,
as they should be for rate constants. A parameter along a direction in which the data give no information is
reported as not identifiable, with an infinite interval. The result also holds the covariance and correlation
matrices; \code{to\_dataframe()} gives a table.

\section{Performance}

Every part of a fit runs inside the compiled engine: simulations, sensitivities, the objective and the
optimizers. Conditions, starts and population members are spread over all cores by one work-stealing thread
pool. A problem needing more than 8\,GB of working memory for its sensitivities is refused when it is built,
with the size it would need; such problems need adjoint sensitivities, which come later.
